% Work-MAT | 2. rocnik | pripravovana hodina 28
% Návaznost: uživatelské připravy.zip (hodiny 01–20), autorská učebnice na GitHubu.
% Realisticky.cz: didaktická inspirace v posloupnosti a typech úloh; zadání jsou vlastní.
\documentclass[11pt,a4paper]{article}
\input{../../sablony/pripravy/priprava.sty}
\nastavPripravu{2. ročník}{Exponenciální nerovnice}{28}{Exponenciální nerovnice I}
\begin{document}
\blok{Motivace a návaznost}
Vrať se ke grafům $2^x$ a $(\frac{1}{2})^x$. Nejdřív nepočítej: žáci mají říci, u které funkce s rostoucím $x$ hodnota roste a u které klesá.
\poznamka{Orientační nabídka pro 45--55 minut: krátké zopakování a motivace, nový výklad, tři společné vzorové úlohy, samostatná práce. Vyber podle tempa třídy; zbývající zadání slouží jako rezerva.}
\blok{Výklad}
Pro $a>1$ platí $a^u>a^v\iff u>v$. Pro $0<a<1$ platí $a^u>a^v\iff u<v$.
\blok{Společné příklady}
\begin{prikladReseniBox}[title={Příklad 1}]
Vyřeš, resp. urči: $2^{2x-1}>2^{x+2}$.
\tcblower
\[\begin{aligned}2x-1>x+2 \\ x>3\end{aligned}\]
\end{prikladReseniBox}
\begin{prikladReseniBox}[title={Příklad 2}]
Vyřeš, resp. urči: $\left(\frac{1}{3}\right)^{x+1}\ge \left(\frac{1}{3}\right)^3$.
\tcblower
\[\begin{aligned}x+1\le3 \\ x\le2\end{aligned}\]
\end{prikladReseniBox}
\begin{prikladReseniBox}[title={Příklad 3}]
Vyřeš, resp. urči: $5^{3x}<25^{x+1}$.
\tcblower
\[\begin{aligned}5^{3x}<5^{2x+2} \\ 3x<2x+2\;\Longrightarrow\;x<2\end{aligned}\]
\end{prikladReseniBox}
\poznamka{Řešení na tabuli zaznačuj na číselné ose. Nepoužívej logaritmy, zatím nejsou probírané.}
\blok{Samostatné procvičování}
\begin{zadaniReseniBox}[title={Samostatná práce 1}]
\begin{enumerate}[label=\alph*)]
\item $2^{x+1}>2^3$.
\item $3^{2x}\le3^{x+3}$.
\item $(\frac{1}{2})^x>(\frac{1}{2})^2$.
\item $5^{x-2}\ge5^0$.
\item $(\frac{1}{3})^{2x}<(\frac{1}{3})^4$.
\item $4^x>2^{x+2}$.
\end{enumerate}
\tcblower
\begin{enumerate}[label=\alph*)]
\item $x>2$.
\item $x\le3$.
\item $x<2$.
\item $x\ge2$.
\item $x>2$.
\item $x>2$.
\end{enumerate}
\end{zadaniReseniBox}
\blok{Dobrovolné rozšíření}
\blok{Úkol na doma}
Při početních úlohách napiš postup, zvolenou metodu a zkoušku v původním zadání. U soustav kontroluj obě rovnice; u nerovnic výsledek zakresli na osu.
\begin{zadaniReseniBox}[title={Domácí úkol}]
\begin{enumerate}[label=\alph*)]
\item $2^{x+2}>8$.
\item $(\frac{1}{2})^{x-1}\ge(\frac{1}{2})^3$.
\item $9^x\le3^{x+3}$.
\item $3^{2x+1}>3^{x-2}$.
\item $(\frac{1}{3})^{x+2}<(\frac{1}{3})^0$.
\item $2^{2x}\ge4^{x-1}$.
\item Kdy obracíme znak nerovnosti?.
\end{enumerate}
\tcblower
\begin{enumerate}[label=\alph*)]
\item $x>1$.
\item $x\le4$.
\item $x\le3$.
\item $x>-3$.
\item $x>-2$.
\item platí pro všechna x.
\item Při základu mezi 0 a 1..
\end{enumerate}
\end{zadaniReseniBox}
\poznamka{Domácí úkol obsahuje 7 položek; předpokládaná domácí práce přibližně 15--30 minut. Po větších blocích jsou vědomě promíchané dříve probrané typy.}
\end{document}
